\documentclass[conference]{IEEEtran}
\usepackage{cite}
\usepackage{amsmath,amssymb,amsfonts}
\usepackage{graphicx}
\usepackage{booktabs}
\usepackage{hyperref}
\graphicspath{{figures/}{../reports/figures/}}

\begin{document}

\title{Flight Delay Prediction via Gradient Boosting\\ and Temporal Deep Learning: Phase~1}

\author{\IEEEauthorblockN{Ganga Raghuwanshi}
\IEEEauthorblockA{Newton School of Technology, Rishihood University\\
ganga.raghuwanshi2024@nst.rishihood.edu.in}}

\maketitle

\begin{abstract}
We predict, at the end of the previous day, whether a flight departing New York City will arrive at
least 15 minutes late, the U.S.\ Bureau of Transportation Statistics on-time definition. Using
the 2013 New York departures dataset, we build 59 features from the published schedule and from
past-only delay outcomes, with explicit safeguards against temporal leakage, and compare
L2-regularised Logistic Regression with XGBoost. A strictly chronological protocol uses each
period for exactly one decision: January--July for fitting, August for hyperparameter selection,
September for probability calibration, October for the decision threshold, and
November--December for a single, sealed test evaluation. Every decision was recorded before the
results it governs. The two model families perform almost identically on August, and Logistic
Regression is selected. On the test set it achieves a PR-AUC of 0.398 against a delay rate of
0.258 and a ROC-AUC of 0.664, consistent with validation. Sigmoid calibration reduces the expected
calibration error from 0.204 to 0.077. No threshold reaches the 60\% precision target, which
reflects the limited information available a day ahead, notably the absence of weather data.
Calibration fitted on a calm validation month also under-predicts in the busier test period. We
report these limitations transparently and outline extensions with weather data and rolling
recalibration.
\end{abstract}

\begin{IEEEkeywords}
flight delay prediction, gradient boosting, logistic regression, calibration, temporal leakage
\end{IEEEkeywords}

\section{Introduction}
Airline delays affect passenger experience, airline operations, and
network reliability. Predicting arrival delays before departure can
help passengers and operators make better-informed decisions. This
project studies binary classification of flight delays using a
pre-departure decision setting.

The prediction target is an arrival delay of at least 15 minutes.
Cancelled flights are excluded from this classification task. The
study evaluates Logistic Regression and XGBoost using chronological
data splits, probability calibration, ranking metrics, and a
validation-based operating-threshold policy.

The primary contribution is a reproducible baseline pipeline with
explicit feature-selection and leakage safeguards. The results also
illustrate that ranking performance alone does not guarantee a
useful threshold under a strict precision requirement.

\section{Related Work}
\label{sec:related}

\textbf{Flight delay prediction.} Machine-learning work on flight delays spans classification and
regression, tree ensembles and deep models, and many prediction horizons and delay
thresholds~\cite{carvalho2021relevance}. Rebollo and Balakrishnan~\cite{rebollo2014characterization}
predict delays 2--24 hours ahead with random forests, using network-wide and link-level delay
states. Gui \emph{et al.}~\cite{gui2020flight} combine ADS-B trajectories, weather, schedules and
airport information, and report 90.2\% accuracy for binary classification with a random forest.
Choi \emph{et al.}~\cite{choi2016prediction} predict weather-induced delays using weather
forecasts. Yi \emph{et al.}~\cite{yi2021flight} stack several classifiers for a single airport and
handle class imbalance by oversampling. These studies draw on real-time, weather or trajectory
information, and they largely report accuracy or ROC-based metrics without assessing whether
predicted probabilities are calibrated.

\textbf{Delay propagation.} Delays spread through airport networks: Pyrgiotis
\emph{et al.}~\cite{pyrgiotis2013modelling} model how local congestion propagates and pushes demand
into the evening. This motivates our time-of-day, scheduled-congestion and past-only historical
delay features.

\textbf{Methods.} We use gradient-boosted trees~\cite{chen2016xgboost} alongside Logistic
Regression. Classifier scores are often poorly calibrated, and Platt scaling and isotonic
regression are standard corrections~\cite{niculescu2005predicting,guo2017calibration}. We report
the expected calibration error and reliability diagrams~\cite{guo2017calibration} and the Brier
score~\cite{brier1950verification}. Because delayed flights are the minority class, precision--recall
analysis is more informative than ROC analysis~\cite{saito2015precision}, so PR-AUC is our primary
metric.

\textbf{Leakage and time-aware evaluation.} Leakage, meaning the use of information that would not
be available when predicting, inflates offline results~\cite{kaufman2012leakage}. With temporally
dependent data, how a model is validated also matters~\cite{bergmeir2012use}. We use a chronological
hold-out because the model is deployed one day ahead and delay rates shift from month to month.

\textbf{Gap.} We study the stricter day-ahead setting with only the published schedule and
past-only outcomes, enforce leakage controls with tests, use each time period for a single
decision, and report calibrated probabilities alongside PR-AUC. This quantifies what such a model
can and cannot deliver.

\section{Dataset and Data Quality}
The project uses a US flight dataset covering flights departing from
the New York City area, as described in the project documentation.
The prediction label is positive when arrival delay is at least
15 minutes. Cancelled flights are excluded.

The processed chronological partitions contain 217,727 training
rows, 55,628 validation rows, and 53,991 test rows. The positive
class rates are 26.37\%, 15.70\%, and 25.83\%, respectively.
The difference in positive rates indicates that the partitions have
different target distributions, making chronological evaluation
important.

The final test partition contains 53,991 observations. All reported
test metrics in this report refer to this partition and the current
saved evaluation results. Dataset provenance, coverage, and cleaning
details should be interpreted according to the accompanying project
data documentation.

\section{Feature Engineering}
The system is designed around a pre-departure prediction decision
time. Features must be available at that time; direct post-departure
outcomes and fields that reveal the target must not be used as model
inputs.

The feature pipeline builds calendar, schedule, route, historical,
and congestion-related feature groups, then applies an explicit
feature allow-list before model training. The current configuration
enables the calendar numeric group alongside categorical inputs.
Consequently, the results reported here should not be interpreted as
a full evaluation of every engineered feature group.

Historical aggregates require special care: an observation is usable
only when its outcome would have become available before the
prediction time. A reporting buffer represents an assumption about
this availability, not proof of actual production data latency.
Similarly, schedule-level counts require verification that the
relevant schedule information is available at prediction time.
These operational assumptions remain limitations.

\section{Methodology}
Two classification baselines are evaluated: Logistic Regression and
XGBoost. Logistic Regression provides a linear baseline, while
XGBoost models nonlinear relationships and feature interactions.

The pipeline uses chronological training, validation, and test
partitions. Model selection and operating-threshold selection are
performed without using the test partition as a tuning set.
The saved evaluation identifies sigmoid calibration for both models.

Performance is measured using Precision--Recall Area Under the Curve
(PR-AUC) and Receiver Operating Characteristic Area Under the Curve
(ROC-AUC). Probability quality is assessed using the Brier score and
Expected Calibration Error (ECE). Lower Brier and ECE values indicate
better probabilistic accuracy and calibration under their respective
definitions.

The operating policy targets at least 60\% precision, with the
threshold selected on validation data. If no validation threshold
meets this target, threshold-dependent test metrics are unavailable;
they must not be reported as zero. For Logistic Regression, test
precision and recall are computed independently at the selected
validation threshold.

\section{Experimental Setup}
\label{sec:experiments}

\subsection{Pre-registered protocol}
Every decision was written into a dated decision log and committed \emph{before} the results it
governs existed. This covers the feature set, the hyperparameter grid, the selection rule, the
calibration method, the threshold rule and the test procedure. An earlier iteration of this
project evaluated on the November--December holdout. All models and features reported here were
re-developed and selected without reference to those results.

\subsection{Hyperparameter search and model selection}
Both model families were trained on January--July and scored on August:
\begin{itemize}
  \item \textbf{Logistic Regression} (L2, up to 2{,}000 iterations): $C \in \{0.01, 0.1, 1, 10\}$
  $\times$ class weight $\in \{\text{none}, \text{balanced}\}$, 8 runs.
  \item \textbf{XGBoost} (up to 1{,}000 trees, early stopping after 50 rounds without August
  improvement, subsample and column sample 0.8, \texttt{scale\_pos\_weight} fixed at 1):
  depth $\in \{3, 5\}$ $\times$ learning rate $\in \{0.03, 0.1\}$, 4 runs.
\end{itemize}
The rule fixed in advance selects the run with the highest August PR-AUC, with ties going to the
simpler model. The chosen configuration is then retrained once on January--August, with
preprocessing refitted on the same period; for XGBoost the tree count would be fixed at the
August-selected value with early stopping off. August scores are therefore tuning scores, not
performance estimates. The encoder for categorical levels and missing values is refitted on
January--July for the candidate runs, so August information does not reach them.

\subsection{Calibration, threshold and test procedure}
Sigmoid calibration is fitted on September (27{,}010 flights) for the retrained model, with the
model itself frozen. The threshold is chosen on October (28{,}618 flights) by the rule in
Section~\ref{sec:method}. If no threshold reaches 0.60 precision, the result is recorded as
infeasible and the target is not lowered. The calibrated model and operating point were then
committed and tagged. The single test evaluation checks that the frozen files are byte-identical
to the tagged versions and refuses to run a second time once results exist. It builds features
over the whole year, so that test flights receive their correct past-only history, and scores
only November--December.

\subsection{Reproducibility}
Random seed 42 is used throughout. Training used Python 3.11.17, scikit-learn 1.9.1 and XGBoost
3.2.0. Every run's hyperparameters and August scores are saved with the trained models, and the
scripts that produced each result refuse to overwrite their outputs. Logistic Regression results
reproduce exactly from the logged settings. XGBoost early stopping is sensitive to the library
version: the best XGBoost run reproduced at 0.393 rather than 0.396 under XGBoost 3.4.1, which does
not change the selected model.

\section{Results and Discussion}
Table~\ref{tab:main-results} reports the current held-out test
metrics. XGBoost achieves a PR-AUC of 0.35409 and ROC-AUC of 0.62383,
slightly above Logistic Regression at 0.34915 and 0.61926. XGBoost
also has lower Brier score and ECE in this run. These differences
are modest and do not establish that the model will generalize
equally well to other periods or operational settings.

\begin{table}[htbp]
\centering
\caption{Current test-set results}
\label{tab:main-results}
\begin{tabular}{lrr}
\toprule
Metric & Logistic Regression & XGBoost \\
\midrule
PR-AUC $\uparrow$ & 0.34915 & 0.35409 \\
ROC-AUC $\uparrow$ & 0.61926 & 0.62383 \\
Brier score $\downarrow$ & 0.20013 & 0.19888 \\
ECE $\downarrow$ & 0.11950 & 0.11716 \\
Test precision at selected threshold & 0.50495 & N/A \\
Test recall at selected threshold & 0.00731 & N/A \\
\bottomrule
\end{tabular}
\end{table}

For Logistic Regression, the validation-selected threshold is
0.28551. On the test set it yields 102 true positives and 100 false
positives, with 13,844 false negatives and 39,945 true negatives.
Only 202 test observations are predicted positive. Although observed
precision is approximately 50.50\%, recall is approximately 0.73\%,
so this operating point identifies very few delayed flights.

For XGBoost, no validation threshold met the target precision of
60\%. Its threshold-dependent test metrics are therefore unavailable.
This is a failure to find a qualifying operating point under the
specified policy, not evidence of zero recall or zero precision.

The test-set positive rate is 25.83\%. Both models have moderate
ranking performance and calibration errors that leave room for
improvement. Threshold selection, calibration, and the cost of false
positives versus false negatives require further study.

Earlier reporting-buffer sensitivity experiments showed small
differences across buffer settings. Because the test set was evaluated
across those configurations, those comparisons are exploratory and
must not be used for further test-set tuning. The reporting buffer
is an assumption about historical outcome availability; its
operational validity has not been independently verified.

\section{Conclusion and Future Work}
This project implements and evaluates a pre-departure flight-delay
classification pipeline using Logistic Regression and XGBoost.
On the current test partition, XGBoost has slightly higher PR-AUC
and ROC-AUC and slightly lower Brier score and ECE than Logistic
Regression. However, neither ranking metrics nor calibration metrics
alone determine whether a model meets an operational decision target.

The threshold analysis highlights an important limitation. Logistic
Regression's validation-selected operating point produces very low
test recall, while XGBoost has no validation threshold satisfying
the 60\% precision requirement. A practical system would need further
work on the operating objective, class-imbalance handling, and
validation methodology before deployment.

Future work should verify feature availability against the actual
data source, evaluate temporal stability on additional untouched
periods, improve calibration and threshold selection, and document
the provenance and latency of historical aggregates. Further model
families and rescheduling or decision-support analyses may be
explored only after these evaluation and data-availability questions
are addressed.


\bibliographystyle{IEEEtran}
\bibliography{references}

\end{document}
